\scanpage{300}
\deleted{\unsure{ce qui exige}} \inserted{\unsure{d'où :}} \deleted{$\unsure{f_1=f_2}$} \inserted{$\unsure{g_1}(t)=\unsure{g_2}(1/t)$} et \uline{$\unsure{g_i}$ est une constante}. Donc les \uline{\unsure{pts} de $Y$ ne sont pas séparés}

On peut aussi \unsure{prendre} un quotient par le groupe $\mathbb Z/2$ [agissant sur $X_1\amalg X_2$ par \unsure{permutation}, et sur $Y$ par $t\mathbin{\unsure{-|\to}}1/t$], on trouve
\margin{\unsure{en réalité ?} \uline{à vérifier}}
\figure{\begin{tikzpicture}\draw (0,.3) .. controls (.5,0) and (.7,.6) .. (1,.3) .. controls (1.5,0) and (1.7,.4) .. (2,.1) .. controls (2.5,-.1) and (3,.15) .. (4,0);\end{tikzpicture}}
\diagram{\begin{tikzcd}
(X_1\amalg_Y X_2)/\mathbb Z_2\\
X_1\amalg X_2\,/\mathbb Z/2\simeq X_1.\arrow[u]
\end{tikzcd}}
\unsure{où|si} $X_1$ est quasi-affine, \deleted{\unsure{comme}} mais non $X_1\amalg_Y X_2\,/\mathbb Z_2$ puisque $X_1\amalg_Y X_2$ le serait \unsure{sinon}.

Il faut cependant vérifier que \unsure{tte} \deleted{partie} \inserted{\unsure{trajectoire}} \deleted{\unsure{trajectoire}} de $X_1\amalg_Y X_2$ par $\mathbb Z/2$ est contenue dans un ouvert affine : c'est facile.
[prendre un \unsure{div} $D$ de $\unsure{\mathfrak X|\mathcal X}$, \inserted{\unsure{contenant le zéro}} \unsure{ne contenant pas} une partie finie, regarder $(\unsure{\mathfrak X|\mathcal X}-D)\amalg(\unsure{\mathfrak X|\mathcal X}-D)$.]
\figure{\begin{tikzpicture}\draw (0,0)--(0,0.4)..controls (0,0.55) and (0.2,0.6)..(0.4,0.58)--(2.5,0.55)..controls (2.6,0.55) and (2.7,0.57)..(2.65,0.68);\end{tikzpicture}}
\mspage{299}
